\section{Mercy Received and Mercy Abused}

St.~Alphonsus describes three acts of divine mercy: God waits for the sinner, calls the sinner, and receives the sinner who returns (\emph{Preparation for Death}, XVI.1--3). Each act concerns a different need. The sinner needs time in which to repent, an awakening from his refusal, and reconciliation when he returns. His conversion depends upon a goodness already at work during his estrangement. Alphonsus asks him to consider that prior mercy so that he can recognize why a present return is possible and give thanks for the help he has received.

The first act is God's patience. Alphonsus asks the sinner to consider the life preserved and the benefits received even during a time of infidelity. The ordinary continuance of life can conceal this gift. We wake, eat, work, make plans, and assume another day will follow. But the time in which a person can repent is not self-created or guaranteed. To find oneself still able to turn toward God is a reason for gratitude.

The second act is the call. Alphonsus names inspirations, the reproach of conscience, preaching, trials, and the deaths of friends among the ways a sinner is awakened. Their common purpose in his meditation is reconciliation. We may have resisted such a summons because it disturbed an arrangement we wanted to preserve. A serious recollection asks whether the same truth has been reaching us repeatedly while we have kept assigning obedience to a later time.

The third act is reception. Alphonsus dwells on Christ's parables of the lost sheep and the returning son. In Luke 15 the father sees the son coming, runs, embraces him, and restores him to the household (Lk 15:11--24). The son has acknowledged his sin and turned homeward. The father's generous reception answers the fear that this return can bring only rejection. The son comes with a truthful confession, and he receives more than he had dared to request. A sinner who has resolved to return can approach Christ's mercy with that confidence.

Shame can become an obstacle at precisely this point. A person may be willing to consider hell in the abstract yet unwilling to say a humiliating truth in confession. St.~Francis de Sales addresses that difficulty after his early meditations on the last things. He urges a courageous, humble confession and distinguishes the shame of committing sin from the goodness of confessing it. He places the penitent, in imagination, at the feet of the crucified Savior, receiving the fruit of his Passion (\emph{Introduction to the Devout Life}, I.19).

Francis places self-accusation before the Savior who suffered for the penitent. Confession brings the sin to Christ for healing through the ministry of the priest. The penitent must therefore speak truthfully even when the admission is humiliating. Concealment would preserve the fault at the moment he has come to seek forgiveness for it. He can accept that humiliation because Christ's Passion has already shown how much his reconciliation matters to the Lord. The counsel and penance he receives then concern the amendment of a life restored by mercy.

Contrition includes sorrow for sin and a purpose of amendment. Fear of eternal punishment can give rise to imperfect contrition, itself a gift of grace, which disposes the sinner for forgiveness in Penance but does not alone remit grave sin. Perfect contrition arises from loving God above all things and obtains forgiveness of mortal sin when joined to the firm resolution of sacramental confession as soon as possible (CCC 1451--1453).

A purpose of amendment is a present willingness to forsake the sin. It does not mean possessing a private certainty that no weakness will ever recur. The penitent can know his frailty while sincerely resolving to resist, use the necessary helps, and return promptly if he falls. But he cannot preserve a deliberate intention to continue the grave evil and call that intention repentance.

Alphonsus addresses the opposite distortion in the next consideration, on the abuse of divine mercy. A person proposes to commit the sin now and confess it afterward. The future act of repentance is used to excuse the present act of rebellion (XVII.1). This is presumption. It treats the mercy which calls us out of sin as though it had promised permission to remain in it until we choose a convenient moment to stop.

The disorder lies in the present choice even before uncertainty about the future is considered. A good intention for tomorrow cannot turn today's deliberate grave evil into obedience. Moreover, tomorrow's repentance is itself a grace; it cannot be put under our control as an item already secured. Presumption claims the benefit of conversion while rejecting conversion's demand. The sinner who has made this bargain can repent of it now, including the misuse of mercy among the sins brought to Christ.

Repentance also reaches what the sin has done to another person. Forgiveness does not make a stolen possession rightfully ours or restore a damaged reputation without our cooperation. We must repair harm as far as possible: return what was taken, correct the false account, make compensation where it is due, and accept the penance given (CCC 1459--1460). Such acts give bodily expression to the change of heart.

The forgiven person must accept that his restored life is a gift. Repeated self-accusation cannot take the place of the amendment and reparation now required. If he has damaged a reputation, continuing to dwell on his guilt while leaving the false account uncorrected prolongs the injury. Receiving mercy should make him more willing to perform the difficult act that benefits his neighbor. The remembrance of hell then includes gratitude for deliverance: he can act as someone whom Christ has received again, using the grace restored to practice the charity his sin contradicted.

